% partes/parte1/cap05.tex
% Capítulo complementario: Google Analytics 4. No corresponde a un subtema del
% programa oficial, por lo que no lleva \progtag (no altera docs/cobertura.md).
\chapter{Google Analytics 4: observar el comportamiento en el sitio propio}
\label{cap:05}

\begin{cajaFicha}
\begin{description}
  \item[Unidad temática] I — Comportamiento, perfil y gestión de datos del consumidor digital (capítulo complementario).
  \item[Temas del programa] Complementa 1.3 (gestión de datos de primera parte) y 1.4.1 (plataformas de gestión integral). No introduce un subtema nuevo del programa oficial.
  \item[Horas] Sesión adicional de 60 a 90 minutos, sin carga de programación.
  \item[Competencia de la unidad] Identifica el comportamiento del consumidor digital y su evolución a partir de sus tipos de perfil y gestión de datos.
  \item[Resultados de aprendizaje observables] Al concluir el capítulo, el estudiante:
  \begin{enumerate}
    \item explica qué es Google Analytics 4 (GA4) y por qué mide eventos y no sesiones;
    \item relaciona una pregunta de mercadotecnia con el informe de GA4 que la responde;
    \item lee un informe de adquisición, un embudo y una serie semanal en la interfaz, sin escribir código;
    \item identifica qué datos de GA4 pueden exportarse, cuándo y con qué precauciones, y las obligaciones de privacidad que implica usarlo.
  \end{enumerate}
  \item[Prerrequisitos] Cap. 1--4.
\end{description}
\end{cajaFicha}

\section{Caso de apertura}

\textbf{Vigente a octubre de 2026.} Quien busque hoy un tutorial de
\enquote{Google Analytics} puede encontrar uno escrito para una herramienta
que ya no existe. La versión anterior, Universal Analytics, dejó de procesar
datos nuevos el 1 de julio de 2023, y el 1 de julio de 2024 se perdió el
acceso a su interfaz y a su API \parencite{google2026ga4ua}. Sin embargo,
buena parte del material en línea, incluidas capturas de pantalla y nombres
de menús, sigue describiendo aquella generación. Es el mismo problema que
abrió el Cap.~4: la plataforma cambia más rápido que el material que la
explica. Por eso este capítulo no enseña una secuencia de clics que
envejecerá, sino el \emph{modelo} con el que GA4 organiza la información
—eventos, parámetros, dimensiones y métricas— y las preguntas de negocio que
cada informe responde. Los nombres de los menús se dan en inglés, tal como
aparecen en la documentación de Google; la interfaz en español puede
traducirlos de otra manera y debe compararse siempre con la ayuda vigente.

\section{Desarrollo teórico}

\subsection{Qué es y de dónde viene}

Google Analytics es un servicio gratuito de analítica de sitios web y
aplicaciones: registra las interacciones de las personas con un sitio o una
app y las resume en informes (existe una versión de pago, Analytics 360, para
organizaciones grandes). Para la mercadotecnia digital es un ejemplo directo
de los \enquote{datos de primera parte} del Cap.~3: información que la
organización recolecta en su propio sitio, a diferencia de la que compra o
recibe de terceros. El cuadro~\ref{tab:cap05-ua-ga4} compara sus dos
generaciones.

\begin{table}[!htb]
\centering
\begin{tabular}{@{}p{3.2cm}p{4.5cm}p{4.7cm}@{}}
\toprule
 & Universal Analytics (UA) & Google Analytics 4 (GA4) \\
\midrule
Unidad de medida & Sesión y vista de página & Evento, con parámetros \\
Estado & Retirada en 2023--2024 & Vigente; sustituyó a UA \\
Plataformas & Sitios web & Sitios web y apps en una misma propiedad \\
\bottomrule
\end{tabular}
\caption{Universal Analytics frente a GA4. Fuente: elaboración propia con base
en la ayuda de Google Analytics \parencite{google2026ga4ua}.}
\label{tab:cap05-ua-ga4}
\end{table}

El cambio de sesiones a eventos no es solo técnico: obliga a decidir
\emph{qué comportamiento se mide} (¿un clic?, ¿ver un video?, ¿completar una
compra?). Es la misma idea del Cap.~2: el perfil de una persona depende de
qué comportamiento se observa.

\subsection{El modelo de datos: eventos, parámetros y propiedades de usuario}

En GA4 todo es un \textbf{evento}: una interacción con un nombre (por
ejemplo, \texttt{page\_view} o \texttt{purchase}) y un conjunto de
\textbf{parámetros} que la describen (la página, el valor de la compra, la
moneda). Además pueden asociarse \textbf{propiedades de usuario} (atributos de
la persona, como un segmento) y, en el comercio electrónico, parámetros de
artículo (producto, categoría, precio). Los eventos se agrupan en cuatro
categorías (cuadro~\ref{tab:cap05-eventos}).

\begin{table}[!htb]
\centering
\begin{tabular}{@{}p{3.1cm}p{5.0cm}p{4.3cm}@{}}
\toprule
Categoría & Qué es & Ejemplos \\
\midrule
Automáticos & Se registran al instalar la etiqueta o el SDK, sin configurar nada & \texttt{first\_visit}, \texttt{session\_start} \\
Medición mejorada & Se activan por defecto en sitios web & \texttt{scroll}, \texttt{click}, \texttt{file\_download} \\
Recomendados & Nombres que Google sugiere y que hay que implementar & \texttt{purchase}, \texttt{add\_to\_cart} \\
Personalizados & Los define la organización & \texttt{compra\_en\_tienda} \\
\bottomrule
\end{tabular}
\caption{Categorías de eventos de GA4. Fuente: elaboración propia con base en
la ayuda de Google \parencite{google2026ga4eventos}.}
\label{tab:cap05-eventos}
\end{table}

Los informes combinan dos tipos de campos. Una \textbf{dimensión} es un
atributo que sirve para agrupar (país, canal de adquisición, página,
dispositivo); una \textbf{métrica} es una cantidad que se mide (usuarios
activos, sesiones, eventos, ingresos). Un informe es, en esencia, una tabla:
una columna por dimensión y otra por métrica.

\begin{cajaAtencion}
\textbf{Límites de recolección.} Para propiedades estándar, la ayuda de Google
indica, entre otros, hasta 25 parámetros por evento, 25 propiedades de usuario
y nombres de evento de 40 caracteres como máximo
\parencite{google2026ga4limites}. Son cifras de una página de ayuda que
cambia; verifique la versión vigente antes de diseñar una medición.
\end{cajaAtencion}

\subsection{De dónde salen los datos y cómo se reconoce a una persona}

Cada fuente de datos es un \textbf{flujo de datos} de la propiedad: un sitio
web con la etiqueta de Google, una app móvil con el SDK de Firebase, o un
servidor que envía eventos mediante el \textit{Measurement Protocol}
(figura~\ref{fig:cap05-flujo}).

\begin{figure}[!htb]
  \centering
  \resizebox{0.94\textwidth}{!}{% figuras/tikz/cap05-flujo-ga4.tex
% Cómo llegan los datos a una propiedad de GA4 y cómo se consultan (vigente a octubre de 2026).
\begin{tikzpicture}[
    nodo/.style={rectangle, rounded corners, draw=guindaIPN, thick, fill=guindaIPNClara,
                 minimum width=3.1cm, minimum height=0.9cm, align=center, font=\footnotesize},
    nodom/.style={rectangle, rounded corners, draw=doradoIPN, very thick, fill=bronceIPNClaro,
                  minimum width=3.3cm, minimum height=2.9cm, align=center, font=\small},
    flecha/.style={-{Stealth[length=2.6mm]}, thick, draw=grisIPN}
  ]
  \node[nodo] (w) at (0,1.4)  {Sitio web\\ (etiqueta de Google)};
  \node[nodo] (a) at (0,0)    {App móvil\\ (SDK de Firebase)};
  \node[nodo] (s) at (0,-1.4) {Servidor o tienda\\ (Measurement Protocol)};
  \node[nodom] (p) at (5,0)   {\textbf{Propiedad GA4}\\[3pt]\footnotesize flujos de datos\\ eventos y parámetros\\ propiedades de usuario};
  \node[nodo] (r) at (10,0.8) {Informes\\ y Exploraciones};
  \node[nodo] (e) at (10,-0.8) {Exportación\\ (solo cuenta propia)};
  \draw[flecha] (w) -- (p); \draw[flecha] (a) -- (p); \draw[flecha] (s) -- (p);
  \draw[flecha] (p) -- (r); \draw[flecha] (p) -- (e);
\end{tikzpicture}
}
  \caption{Cómo llegan los datos a una propiedad de GA4 y cómo se consultan.
  La exportación solo está disponible en una cuenta propia (ver §4). Fuente:
  elaboración propia.}
  \label{fig:cap05-flujo}
\end{figure}

GA4 no \enquote{conoce} a nadie: reconoce \emph{navegadores y dispositivos}
mediante un identificador de cliente guardado en la cookie \texttt{\_ga}, que
por defecto caduca a los 2 años \parencite{google2026ga4cookies}. Si la
organización tiene inicios de sesión, puede enviar además un identificador
propio del usuario (\textit{User-ID}). Una persona con tres dispositivos y sin
sesión iniciada aparece como tres usuarios: es el problema de la identidad
única de la omnicanalidad (Cap.~3).

\subsection{Para qué sirve en mercadotecnia digital}

El cuadro~\ref{tab:cap05-usos} relaciona preguntas de negocio con el informe
de GA4 que las responde y con el capítulo del libro donde se estudió el
concepto.

\begin{cajaFrontera}
\textbf{Un dato es tan bueno como la pregunta.} GA4 mide comportamiento
\emph{agregado y anónimo} dentro de \emph{su} sitio o app: no dice por qué la
gente actúa así ni qué hace fuera de él. Complementa a la encuesta, la escucha
social y el CRM.
\end{cajaFrontera}

\begin{table}[!htb]
\centering
\small
\begin{tabular}{@{}p{4.2cm}p{5.4cm}p{2.8cm}@{}}
\toprule
Pregunta de negocio & Dónde se responde & Se conecta con \\
\midrule
¿De dónde llegan mis visitantes y qué canales convierten? & \textit{Acquisition} (adquisición de usuarios y de tráfico) & Cap.~4 (atribución) \\
¿Qué hacen y en qué paso abandonan? & \textit{Engagement}; embudo y rutas & Cap.~1 (fricción) \\
¿Qué se vende y cuánto ingreso deja? & \textit{Monetization} & Cap.~3 (valor) \\
¿Vuelven? ¿Cuánto vale una cohorte? & \textit{Retention}; cohortes y valor de vida & Cap.~3 (CLV) \\
¿Qué segmentos tengo? & Audiencias; superposición de segmentos & Cap.~2 \\
¿Hay estacionalidad en mi tráfico? & Serie semanal leída y descompuesta & Cap.~4 \\
\bottomrule
\end{tabular}
\caption{Preguntas de negocio y su informe en GA4. Fuente: elaboración propia.}
\label{tab:cap05-usos}
\end{table}


\section{Evidencia sobre cómo visualizar los datos}

\subsection{La interfaz y los informes estándar}

La interfaz se organiza en cinco áreas: \textit{Home} (resumen),
\textit{Reports} (informes estándar), \textit{Explore} (análisis a la medida),
\textit{Advertising} (atribución y planeación de canales) y \textit{Admin}
(configuración: flujos, eventos clave, retención). Los informes estándar se
agrupan en colecciones; la predeterminada, \textit{Life cycle}, sigue el
recorrido del cliente: \textit{Acquisition}, \textit{Engagement},
\textit{Monetization} y \textit{Retention} \parencite{google2026ga4informes}.
Las acciones que importan al negocio (compra, registro) se marcan como
\textbf{eventos clave}, término que Google usa en Analytics y que reserva
\enquote{conversión} para el contexto de anuncios \parencite{google2026ga4eventosclave}.

\subsection{Exploraciones}

\textit{Explore} permite construir análisis propios con siete técnicas
\parencite{google2026ga4explorar}: \textit{Free form} (tablas cruzadas y
gráficas), \textit{Funnel} (embudo), \textit{Path} (rutas), \textit{Segment
overlap} (traslape de segmentos), \textit{User explorer}, \textit{Cohort} y
\textit{User lifetime}. De ellas, este capítulo usa dos: el embudo, que
muestra en qué paso de una tarea se pierden los usuarios (la fricción del
Cap.~1), y la tabla libre con gráfica de líneas, para leer una serie semanal.

\begin{cajaAtencion}
\textbf{Retención de datos.} En propiedades estándar, los datos de usuario y
de evento se conservan 2 o 14 meses, según se configure. Ese ajuste limita las
Exploraciones y los embudos, \emph{no} los informes estándar agregados
\parencite{google2026ga4retencion}. Una exploración de más de 14 meses se
\enquote{corta} aunque el informe estándar sí muestre ese periodo; y para
distinguir tendencia de estacionalidad (Cap.~4) hacen falta al menos dos
ciclos anuales.
\end{cajaAtencion}

\section{Laboratorio en pantalla: la cuenta de demostración}

\begin{cajaLaboratorio}[title={Laboratorio en pantalla (sin código)}]
\textbf{Objetivo.} Leer en la interfaz de GA4, sin escribir ni ejecutar código,
tres piezas de información: de dónde llega el tráfico, dónde abandonan los
usuarios y qué forma tiene una serie semanal.

\textbf{Cuenta.} Google ofrece una cuenta de demostración con datos reales de
comercio electrónico de su tienda de mercancía (\textit{Google Merchandise
Store}) \parencite{google2026ga4demo}. Se abre desde \url{https://goo.gle/ga-demo}
con una cuenta de Google; en el selector de propiedades se elige
\textbf{GA4 -- Google Merchandise Store}. Es una propiedad de solo consulta
(rol \textit{Viewer}): no se puede cambiar la configuración, y eso es lo
esperado. Se trata de datos de Google, no del curso: conviene confirmar en el
enlace oficial que la propiedad sigue disponible.

\textbf{Aviso: la cuenta de demostración no permite exportar.} En la práctica
con el grupo, la opción \textit{Export data} no estuvo disponible en esta
cuenta: no se pueden descargar ni exportar datos. Es una limitación de la
cuenta pública de demostración, no un error de quien la usa. En una
\textbf{propiedad propia}, con un rol suficiente (por ejemplo, Analista o
superior; verifique en la ayuda vigente qué rol exige cada exportación), sí se
podrá exportar a CSV, Hojas de cálculo o PDF. Por eso las tres tareas
siguientes se resuelven \textbf{leyendo y capturando la pantalla}: ninguna
exige descargar archivos.

\textbf{Tarea 1: ¿de dónde llega el tráfico?} En \textit{Reports}, abra
\textit{Acquisition $\rightarrow$ Traffic acquisition}; fije los últimos 28
días y cambie la dimensión principal a \textit{Session source / medium}.
Identifique el canal con más sesiones y el de mejor tasa de interacción.
¿Coinciden? Explique la diferencia en una frase.

\textbf{Tarea 2: un embudo de compra.} En \textit{Explore}, cree una
exploración \textit{Funnel} con los pasos \texttt{view\_item},
\texttt{add\_to\_cart}, \texttt{begin\_checkout} y \texttt{purchase}. Anote el
porcentaje que pasa de un paso al siguiente y señale el de mayor abandono.
Compare móvil y escritorio.

\textbf{Tarea 3: leer una serie en pantalla.} En \textit{Explore}, cree una
exploración \textit{Free form} con visualización de líneas: dimensión semana
(o fecha), métrica sesiones, el rango más largo disponible. Tome una captura.
Sobre la captura, marque con lápiz la \textbf{tendencia} (¿sube, baja, se
mantiene?) y los \textbf{picos que se repiten} cada año. Responda en una o dos
frases si el patrón es de temporada o de campaña, y si hay dos ciclos
completos para confirmarlo; es exactamente la pregunta del estudio de caso
del Cap.~4.

\textbf{Opcional, solo con cuenta propia.} Exporte la serie a CSV, déjela en
Excel con dos columnas (fecha y valor) y súbala a la aplicación
\texttt{cap04\_series\_datos\_propios.html} del Cap.~4 con 52 periodos por
ciclo, para obtener la descomposición automática.
\end{cajaLaboratorio}

\begin{cajaAtencion}
\textbf{Qué esperar de un dato exportado.} Lo que sale de un informe viene
\emph{agregado} (no hay filas por persona) y puede tener umbrales de
privacidad o muestreo en exploraciones grandes; dos fuentes de la misma
propiedad pueden no coincidir exactamente. Documente siempre de qué fuente,
con qué filtros y con qué fecha de consulta salió cada cifra, el mismo
criterio de verificación de este libro.
\end{cajaAtencion}

\subsection{Otras vías para llevarse los datos}

Además de exportar desde una exploración, GA4 ofrece conectar la propiedad a
Looker Studio (tableros interactivos para reportes periódicos), exportar los
eventos crudos a BigQuery (un almacén de datos en la nube de Google, que
configura el área técnica y puede generar costos) y consultar los datos desde
otros programas mediante una API. Las dos últimas requieren programación o
apoyo técnico y no se usan en este curso; se mencionan para que el estudiante
sepa que existen y a quién pedirlas.

\begin{cajaEtica}
Medir comportamiento con una cookie que identifica un navegador, con un
identificador de usuario o con un CRM unido a la analítica es tratamiento de
datos personales, y el Cap.~3 lo sitúa bajo la LFPDPPP: se necesita base legal
y propósito declarado (verifique el texto vigente de la ley). El sitio debe
informar el uso de cookies y analítica, y GA4 ofrece un \textit{Consent Mode}
que ajusta lo que miden las etiquetas según el consentimiento del usuario
\parencite{google2026ga4consent}. Principio de minimización: recolecte solo lo
necesario, nunca nombres, correos ni teléfonos en un campo de analítica, y
elija la retención de 2 o 14 meses según la necesidad real. La documentación
de Google indica que Analytics no registra ni almacena direcciones IP; aun
así, la organización sigue siendo responsable del resto de lo que recolecta.
\end{cajaEtica}

\section{Estudio de caso}

\begin{casoFicticio}
Una marca ficticia de ropa deportiva vende en línea y en tiendas físicas. Su
equipo abre GA4 y observa que, en los últimos 28 días, el canal con más
sesiones es la búsqueda orgánica, pero el que tiene mejor tasa de interacción
es un correo de lanzamiento. Al armar el embudo, nota que el mayor abandono
ocurre entre \texttt{add\_to\_cart} y \texttt{begin\_checkout}, sobre todo en
móvil. Además, la gráfica semanal de dos años muestra un pico cada noviembre.
La dirección propone duplicar el presupuesto del correo y atribuirle el pico
de noviembre.

\textbf{Preguntas guía.}
\begin{enumerate}
  \item ¿Por qué el canal con más sesiones puede no ser el de mejor tasa de
  interacción, y qué decisión distinta sugiere cada métrica?
  \item ¿Qué fricción (Cap.~1) podría explicar el abandono entre agregar al
  carrito e iniciar el pago en móvil, y qué cambio probaría usted?
  \item ¿Qué información adicional necesita para saber si el pico de
  noviembre es temporada o efecto del correo? Relacione su respuesta con el
  Cap.~4.
  \item La marca vende también en tienda. ¿Qué parte de su cliente real
  \emph{no} ve GA4, y qué otra fuente de datos propia la completaría?
\end{enumerate}
\end{casoFicticio}

\section{Actividades y evidencias del portafolio}

\begin{description}
  \item[Reporte de la práctica en pantalla] Reporte breve (2 a 3 páginas) con
  las capturas de las tres tareas del laboratorio y las respuestas a sus
  preguntas. Rúbrica sugerida (100 puntos): contexto 10, acceso y capturas 10,
  adquisición 20, embudo 25, lectura de la serie y estacionalidad 25,
  reflexión con la matriz del Cap.~4 10.
  \item[Reflexión con la matriz del Cap.~4] Evaluar a Google Analytics con los
  cuatro criterios (modelo de negocio, intermediación de datos,
  interoperabilidad y exposición regulatoria): ¿quién es dueño del dato?, ¿qué
  tan fácil es llevarse los datos a otra herramienta?
  \item[Solución de casos] Resolver por escrito el estudio de caso de la marca
  de ropa deportiva (§5).
\end{description}

\section{Ruta de aprendizaje autónomo}

Tarea concreta: abrir la ayuda vigente de Google Analytics y comprobar si
cambió algo de lo descrito en este capítulo (nombre de menús, rol que exige
exportar, límites de recolección o plazos de retención), anotando la fecha de
consulta y la página. El capítulo se cerró con la ayuda consultada el 12 de
octubre de 2026.

\section{Síntesis}

GA4 es la herramienta con la que una organización observa el comportamiento
de las personas en su propio sitio o app. Su unidad es el evento con
parámetros; sus informes son tablas de dimensiones y métricas; reconoce
navegadores, no personas. Se aprende el modelo y no los clics, porque la
interfaz cambia. Lo que se exporta viene agregado y con restricciones, y la
cuenta de demostración no permite exportar: en una cuenta propia sí.
Recolectar y analizar estos datos es tratamiento de datos personales y exige
base legal, aviso y minimización.

\subsection*{Glosario del capítulo}
\begin{description}
  \item[Dimensión] Atributo para agrupar los datos en un informe (país, canal, página).
  \item[Evento] Interacción con nombre y parámetros (p.\,ej. \texttt{purchase}); unidad de medición de GA4.
  \item[Evento clave] Evento marcado como importante para el negocio (antes, \enquote{conversión} en Analytics).
  \item[Exploración] Análisis a la medida: tabla libre, embudo, rutas, cohortes.
  \item[Flujo de datos] Fuente de datos de una propiedad: sitio web o app.
  \item[Métrica] Cantidad que se mide (sesiones, usuarios activos, ingresos).
  \item[Parámetro] Dato que describe un evento (valor, moneda, página).
  \item[Propiedad de usuario] Atributo de una persona asociado a sus eventos (p.\,ej. un segmento).
\end{description}

\subsection*{Autoevaluación}

\textbf{Opción múltiple}

\begin{enumerate}
  \item La unidad básica de medición de GA4 es:
  \begin{enumerate}
    \item[a)] La sesión.
    \item[b)] El evento, con sus parámetros. \textit{(correcta)}
    \item[c)] La página vista.
    \item[d)] El usuario identificado por correo.
  \end{enumerate}
  \textit{Justificación:} §2.2.

  \item La retención de datos de 14 meses:
  \begin{enumerate}
    \item[a)] Borra los informes estándar anteriores.
    \item[b)] Limita las Exploraciones y los embudos, no los informes estándar. \textit{(correcta)}
    \item[c)] Impide exportar datos.
    \item[d)] Solo aplica a las apps móviles.
  \end{enumerate}
  \textit{Justificación:} §3.2.

  \item Los datos que se exportan de un informe de GA4 vienen:
  \begin{enumerate}
    \item[a)] Con una fila por persona.
    \item[b)] Agregados, por lo que hay que documentar su fuente. \textit{(correcta)}
    \item[c)] Sin ninguna restricción de privacidad.
    \item[d)] Siempre idénticos a los de cualquier otra exportación.
  \end{enumerate}
  \textit{Justificación:} §4.
\end{enumerate}

\textbf{Preguntas abiertas}

\begin{enumerate}
  \item Explique por qué la cuenta de demostración de Google Analytics no
  sirve para practicar la exportación y cómo se resuelve la práctica de este
  capítulo sin descargar datos. \textit{Clave:} es una propiedad de solo
  consulta; la exportación sí estará disponible en una propiedad propia con
  rol suficiente; las tareas se resuelven leyendo y capturando la pantalla
  (§4).
  \item Explique con sus palabras por qué un salto de 20 puntos en las
  últimas semanas no basta para atribuirlo a una campaña. \textit{Clave:}
  debe mencionar la estacionalidad y la necesidad de comparar con el mismo
  periodo del año anterior (dos ciclos), como en el Cap.~4.
\end{enumerate}

\section{Lecturas recomendadas}

\begin{description}
  \item[\textcite{google2026ga4ua}] Aclara qué generación de la herramienta
  está vigente y evita usar tutoriales de Universal Analytics por error.
  \item[\textcite{google2026ga4explorar}] Describe las técnicas de exploración
  con más detalle que este capítulo, útil para ampliar la práctica con una
  cuenta propia.
\end{description}
